\section{Task C: Vary the Duty Cycle}

\subsection{a) Calculations:}
In this Exercise we will implement a Pulse counter with Timer 3 which counts the overflows of Timer 2.

For this we will calculate the Prescaler and ARR values for Timer 2 again the same as in the first two tasks.
\begin{align}
	f_{clk} &= \SI{75}{\mega\hertz}
\end{align}

We will use Formula \ref{eq:timer3} to calculate the Prescaler value (PSC) and the Counter Period (ARR).
\begin{align}\label{eq:timer3}
	f = \frac{f_{clk}}{(PSC + 1)(ARR + 1)}
\end{align}

From this we can solve for ARR to get
\begin{align}\label{eq:arr}
	ARR &= \frac{f_{clk}}{f(PSC + 1)} - 1
\end{align}

With the chosen Prescaler being $5000$, the PSC value results in 
\begin{align}
	PSC + 1 &= 5000 \\
	PSC &= 4999
\end{align}

Now we can calculate our ARR using \ref{eq:arr}, which results in our maximal Counter Value being 
\begin{align}
	ARR &= \frac{\SI{75}{\mega\hertz}}{\SI{1}{\hertz} \cdot \num{5000}} - 1 = \num{14999}
\end{align}

With these values the desired Frequency of $\SI{1}{\hertz}$ is generated in Timer 2.

For Timer 3 we don't need to calculate anything, because this only counts the overflows of Timer 2.

\subsection{b) Implementation:}
Now we can configure the timer in the STM32CubeMX Program as seen in Figure \ref{fig:config_timer2}.

\begin{figure}[H]
	\centering
	\includegraphics[width=0.5\linewidth]{config_timer2.png}
	\caption{STM32CubeMX configuration of Timer 2 with the ARR and PSC values entered}
	\label{fig:config_timer2}
\end{figure}

And the configuration of the counter in the STM32CubeMX Program can be seen in Figure \ref{fig:config_timer3}. Important to note here, is that the Trigger for this counter is the Timer interrupt ITR1 and the clock Source also needs to be changed to the External Clock. This ensures, that the Counter uses timer 2 as a timing source and not the internal clock.

\begin{figure}[H]
	\centering
	\includegraphics[width=0.5\linewidth]{config_timer3.png}
	\caption{STM32CubeMX configuration of the Counter}
	\label{fig:config_timer3}
\end{figure}

The GPIO Pins also have to be configured the same as in Task A.

\pagebreak
Now we will implement a method to set the correct bits on the 7-Segment Display.

\begin{lstlisting}[language=C, numbers=none, caption={Declaration and Iplementation of the 7-Seg Display mapping function}, captionpos=b]
void set7Seg(uint8_t digit)
{
	static const uint32_t digit_pat[10] = {
		0b0111111, // 0
		0b0000110, // 1
		0b1011011, // 2
		0b1001111, // 3
		0b1100110, // 4
		0b1101101, // 5
		0b1111101, // 6
		0b0000111, // 7
		0b1111111, // 8
		0b1101111  // 9
	};
	
	if (digit > 9U)
	{
		return;
	}
	
	/* Mask for the lower 7 bits */
	uint32_t mask = 0b1111111;
	
	/* Clear segments A..G and then set the pattern for the selected digit */
	GPIOJ->ODR &= ~mask;
	GPIOJ->ODR |= (digit_pat[digit] & mask);
}
\end{lstlisting}

Then in the main function we will start the Timers and set the current counter value.

\begin{lstlisting}[language=C, numbers=none, caption={Starting both counters and setting the current counter value}, captionpos=b]
  /* Reset counters to 0 */
__HAL_TIM_SET_COUNTER(&htim2, 0);
__HAL_TIM_SET_COUNTER(&htim3, 0);

/* Start TIM2: 1 Hz as internal trigger */
HAL_TIM_Base_Start(&htim2);

/* Start TIM3: counts external clock from TIM2 */
HAL_TIM_Base_Start(&htim3);

/* Read counter (0..9) and display it on the 7-seg */
uint8_t curr_cnt = (uint8_t)(__HAL_TIM_GET_COUNTER(&htim3) % 10U);
set7Seg(curr_cnt);
\end{lstlisting}

\pagebreak
Then in the while loop we update the 7-Segment display only when the counter value has changed. 

\begin{lstlisting}[language=C, numbers=none, caption={Polling Timer 3 to update the 7-Segment Display}, captionpos=b]
while (1)
{
	/* Read current counter value from TIM3 */
	uint8_t new_cnt = (uint8_t)(__HAL_TIM_GET_COUNTER(&htim3) % 10U);
	
	/* Update display only when the value changes */
	if (new_cnt != curr_cnt)
	{
		curr_cnt = new_cnt;
		set7Seg(curr_cnt);
	}	
}
\end{lstlisting}

\subsection{c) Results:}
Here in Figure \ref{fig:task_c_2} the Counter value is set to 2 and this is correctly displayed on the 7-Segment display.

\begin{figure}[H]
	\centering
	\includegraphics[angle= 90, width=0.5\linewidth]{task_c_2.JPEG}
	\caption{Counter value 2 displayed on the 7-Segment}
	\label{fig:task_c_2}
\end{figure}

And here in Figure \ref{fig:task_c_9} the LED can be seen as it is turned on at full brightness.

\begin{figure}[H]
	\centering
	\includegraphics[angle= 90, width=0.5\linewidth]{task_c_9.JPEG}
	\caption{Counter value 9 displayed on the 7-Segment}
	\label{fig:task_c_9}
\end{figure}

\subsection{d) Discussion:}
In this task, two timers work together to create a simple pulse counter. Timer 2 serves as the master timer, generating a precise one-second time base. Timer 3 acts as the slave, using that signal as its external clock input. With this setup, Timer 3 automatically counts from 0 to 9 without any software controlling the timing.
The application just needs to read Timer 3's counter value and update the seven-segment display whenever the count changes. By using fixed bit patterns and a bit mask for the display (as in Practicum 1, Task D), the code stays compact and easy to follow. Each digit corresponds to a predefined pattern, and all display segments are updated in one write operation. The set7Seg() function works as a reusable snippet for handling this.
